% 17

\eocesol{(a)~Syarat-syarat binomial terpenuhi:
(1)~Percobaan yang saling independen: Dalam sampel acak, apakah seorang berusia
18--20 tahun telah mengonsumsi alkohol atau tidak, tidak bergantung pada
apakah orang lain telah mengonsumsinya atau tidak.
(2)~Jumlah percobaan yang tetap: $n = 10$.
(3)~Hanya ada dua hasil pada setiap percobaan: Mengonsumsi atau tidak
mengonsumsi alkohol.
(4)~Peluang sukses sama untuk setiap percobaan: $p = 0.697$.
(b)~0.203.
(c)~0.203.
(d)~0.167.
(e)~0.997.}

% 19

\eocesol{(a)~$\mu = 35$, $\sigma = 3.24$
(b)~$Z = \frac{45 - 35}{3.24} = 3.09$. Nilai 45 berjarak lebih dari 3
simpangan baku dari rata-rata sehingga kita dapat menganggapnya sebagai
pengamatan yang tidak biasa. Karena itu, ya, kita akan terkejut.
(c)~Dengan menggunakan pendekatan normal terhadap distribusi binomial,
0.0010. Dengan koreksi 0.5,
0.0017.}

% 21

\eocesol{(a)~$1-0.75^3 = 0.5781$.
(b)~0.1406.
(c)~0.4219.
(d)~$1-0.25^3=0.9844$.}

% 23

\eocesol{(a)~Distribusi geometrik: 0.109.
(b)~Binomial: 0.219.
(c)~Binomial: 0.137.
(d)~$1-0.875^6=0.551$.
(e)~Geometrik: 0.084.
(f)~Dengan menggunakan distribusi binomial dengan $n = 6$ dan $p=0.75$,
kita memperoleh $\mu=4.5$, $\sigma=1.06$, dan $Z = 2.36$. Karena nilai ini
tidak berada dalam jarak 2 simpangan baku, pengamatan tersebut dapat dianggap
tidak biasa.}

% 25

\eocesol{(a)~$\stackrel{Anna}{1/5}\times\stackrel{Ben}{1/4}\times\stackrel{Carl}{1/3}\times\stackrel{Damian}{1/2}\times\stackrel{Eddy}{1/1} = 1/5!=1/120$.
(b)~Karena peluang-peluang tersebut harus berjumlah 1, harus ada $5!=120$ kemungkinan
urutan.
(c)~$8!=\text{40,320}$.}
